\label{ch:introduction}

\section{Purpose and how to read this document}

This document is the inference-side companion to the \emph{Theoretical
Foundations of \texttt{tpeanuts}} document: where that document develops the
generic neutrino-propagation physics implemented in \texttt{core/} and
\texttt{medium/}, this one documents a single, complete, real-data
application built on top of it -- fitting the reactor mixing angle
$\theta_{13}$ and the mass-squared splitting $\Delta m^2_{31}$ to the real
public data release of the Daya Bay reactor antineutrino experiment.
Every number quoted here that is not explicitly attributed to a fitted
\texttt{tpeanuts} result is a real, published, externally verifiable
quantity (a detector specification, a reactor parameter, or a Daya Bay
Collaboration measurement), traceable through the References chapter.

The five chapters are organised as follows: this chapter introduces the
experiment itself -- geometry, detector design, and detection principle --
with no reference yet to \texttt{tpeanuts}; \cref{ch:data-release}
catalogues the real official data release the code consumes, table by
table; \cref{ch:detector-implementation} is the physics-and-mathematics
core, developing every formula \texttt{tpeanuts.detector.dayabay}
implements; \cref{ch:inference-process} documents the statistical inference
machinery (likelihoods, fits, likelihood scans) and the three fit
strategies used; \cref{ch:results} presents, interprets, and summarises the
numerical results obtained from \texttt{inference4\_dayabay.ipynb},
compared against Daya Bay's own published values and against NuFIT.

\section{Overview}

The Daya Bay Reactor Neutrino Experiment is a multi-detector, multi-reactor
survey of electron-antineutrino disappearance, located at the Daya Bay
nuclear power complex in Shenzhen, Guangdong province, China. It was
designed with a single, focused physics goal: to measure the smallest of
the three PMNS mixing angles, $\theta_{13}$, via the disappearance channel
\begin{equation}
  \bar\nu_e \to \bar\nu_e, \qquad
  P_{\bar\nu_e\bar\nu_e}(E,L) \simeq 1 - \sin^2 2\theta_{13}\,
    \sin^2\!\left(\frac{\Delta m^2_{ee}L}{4E}\right)
    + \left(\theta_{12},\,\Delta m^2_{21}\ \text{terms}\right),
  \label{eq:intro-survival-schematic}
\end{equation}
using only reactor-produced $\bar\nu_e$ (no beam, no accelerator) and pairs
of identically designed detectors at different baselines. Data taking began
in December 2011 and continued through 2020, accumulating more than 5.55
million inverse-beta-decay (IBD) candidates with neutron capture on
gadolinium (nGd) over 3158 days of operation
\parencite{An2017DayaBayFluxSpectrum}. The experiment reported the first
$>5\sigma$ observation of a non-zero $\theta_{13}$ in March 2012
\parencite{An2012DayaBayDiscovery}, and subsequently refined the
measurement to few-percent precision on both $\sin^2 2\theta_{13}$ and the
effective splitting $\Delta m^2_{ee}$ \parencite{An2023DayaBayPrecision}.
This document, and the \texttt{tpeanuts} code it accompanies, is built
against the Collaboration's own full public data release
\parencite{DayaBayZenodo2025}, the same underlying dataset behind
\cite{An2023DayaBayPrecision}.

\begin{notebox}
\textbf{Scope of this document.} \texttt{tpeanuts} reproduces the real
6-reactor/8-detector geometry, the real reactor antineutrino flux (with its
real non-equilibrium and spent-fuel corrections), the real order-1/M
inverse-beta-decay cross section, and the real detector energy response
(IAV, LSNL, Gaussian resolution) at their nominal/central values. It does
\emph{not} reproduce the full correlated systematic covariance matrix of
the official Daya Bay combined analysis, nor the 6AD/7AD sub-periods (only
the 8AD period, all 8 detectors online, is modelled). \Cref{ch:results}
reports, honestly and without adjustment, where this simplification does
and does not matter for the final $\theta_{13}$/$\Delta m^2_{31}$ result.
\end{notebox}

\section{The reactor complex: six real cores, three nuclear power plants}
\label{sec:reactor-geometry}

\begin{theorybox}
The antineutrino sources are six pressurised-water reactor cores, grouped
in pairs at three nuclear power plants along the Daya Bay/Ling Ao
site \parencite{An2017DayaBayFluxSpectrum,DayaBayZenodo2025}:
\begin{itemize}
  \item \textbf{Daya Bay NPP}: reactor cores D1, D2.
  \item \textbf{Ling Ao NPP}: reactor cores L1, L2.
  \item \textbf{Ling Ao-II NPP}: reactor cores L3, L4.
\end{itemize}
In the official data release and throughout this document (and the
\texttt{tpeanuts} code), the six cores are labelled generically
$\mathrm{R1}$--$\mathrm{R6}$. Each core has a real nominal thermal power of
\begin{equation}
  P_{\rm th} \approx 2.895\ \mathrm{GW_{th}},
\end{equation}
a single shared, time-averaged value in the data release used identically
for all 6 cores (the real per-core power history varies week to week; see
\cref{sec:weekly-rate-table}). Antineutrinos are produced by the beta decays
of neutron-rich fission fragments of four fissile isotopes present in the
reactor fuel, $^{235}$U, $^{238}$U, $^{239}$Pu, $^{241}$Pu, with real
time-averaged fission fractions (shared by all 6 cores in this data tier)
\begin{equation}
  f_{^{235}{\rm U}} \approx 0.563,\quad
  f_{^{238}{\rm U}} \approx 0.076,\quad
  f_{^{239}{\rm Pu}} \approx 0.305,\quad
  f_{^{241}{\rm Pu}} \approx 0.056.
\end{equation}
On average, about six antineutrinos are released per fission, but only
those above the inverse-beta-decay threshold ($E_\nu \gtrsim 1.8$~MeV,
\cref{ch:detector-implementation}) are detectable -- roughly two per
fission, integrated over the real tabulated Huber-Mueller spectra
(\cref{ch:data-release}).
\end{theorybox}

\begin{notebox}
Every one of the 6 real cores is treated as a point source located at its
own real, published position; the finite physical extent of a reactor core
introduces a negligible baseline correction (a few centimetres against
baselines of hundreds of metres) and is not modelled separately -- this
matches the official analysis's own documented treatment
\parencite{An2017DayaBayFluxSpectrum}.
\end{notebox}

\section{The detector array: three experimental halls, eight antineutrino detectors}
\label{sec:detector-geometry}

\begin{theorybox}
Eight functionally identical antineutrino detectors (ADs) are deployed
underground in three experimental halls (EH), chosen to sit at baselines
spanning the full oscillatory pattern of \labelcref{eq:intro-survival-schematic}:
\begin{itemize}
  \item \textbf{EH1} (near hall): 2 ADs (AD11, AD12), real flux-weighted
    baseline $\approx 560$~m.
  \item \textbf{EH2} (near hall): 2 ADs (AD21, AD22), real flux-weighted
    baseline $\approx 600$~m.
  \item \textbf{EH3} (far hall): 4 ADs (AD31, AD32, AD33, AD34), real
    flux-weighted baseline $\approx 1640$~m.
\end{itemize}
``Flux-weighted baseline'' means the effective distance obtained by
weighting each of the 6 real (detector, reactor) baselines by that
reactor's real relative antineutrino output -- since every detector
actually sees a superposition of all 6 cores, not a single one
(\cref{eq:dayabay-flux-sum} in \cref{ch:detector-implementation}). EH1 and
EH2 together form the ``near'' halls (only weakly oscillated, used as the
reference spectrum in a near/far comparison, \cref{ch:inference-process});
EH3 is the ``far'' hall, sitting close to the first oscillation minimum for
the bulk of the reactor antineutrino energy spectrum and hence carrying
most of the real $\theta_{13}$-driven deficit.

The detector-naming convention used throughout this document and the
\texttt{tpeanuts} code, \code{AD}$hd$, encodes the hall number $h$ and a
sequential detector index $d$ within that hall directly in the name (e.g.\
\code{AD31} is the first AD installed at EH3) -- this is also how
\texttt{tpeanuts.detector.dayabay.parameters.NEAR\_DETECTORS}/
\texttt{FAR\_DETECTORS} group detectors by hall
(\cref{ch:detector-implementation}).
\end{theorybox}

\subsection*{Detector construction}

\begin{theorybox}
Each AD is built from three nested cylindrical vessels, designed to be as
close to functionally identical across all 8 detectors as physically
achievable (a deliberate design choice: systematic detector-to-detector
differences would otherwise be indistinguishable from a real oscillation
effect in a near/far comparison) \parencite{An2017DayaBayFluxSpectrum}:
\begin{enumerate}
  \item \textbf{Inner Acrylic Vessel (IAV)}, wall thickness $11$~mm,
    containing $0.1\%$ gadolinium-doped liquid scintillator (GdLS) --
    the actual neutrino target volume. Gadolinium's huge thermal-neutron
    capture cross section ($\sigma\sim 4.9\times10^4$~b for natural Gd,
    versus $\sim0.3$~b for hydrogen) is what makes the delayed-coincidence
    tag efficient and fast (\cref{sec:ibd-principle}).
  \item \textbf{Outer Acrylic Vessel (OAV)}, filled with undoped liquid
    scintillator (LS) -- a ``gamma catcher'' that increases the detection
    efficiency for gamma rays (from both the prompt positron-annihilation
    signal and the delayed neutron-capture signal) that escape the smaller
    IAV volume without being fully absorbed there.
  \item \textbf{Outermost stainless-steel tank}, filled with mineral oil,
    hosting a total of $192$ 8-inch photomultiplier tubes (PMTs) radially
    positioned around the two inner vessels, plus specular reflectors above
    and below the outer acrylic vessel to improve light collection
    uniformity.
\end{enumerate}
Three automated calibration units (ACUs), each capable of deploying
radioactive sources along one of three vertical axes, sit atop each AD's
outer tank, providing the real energy-scale/position calibration data
underlying the LSNL and IAV response corrections
(\cref{sec:iav-lsnl-theory}).

Each AD is submerged in a two-zone water-Cherenkov muon-veto system (inner
and outer water shield, IWS/OWS), instrumented with its own PMTs, providing
the real muon-veto efficiency folded into the real daily
\code{eff\_livetime} exposure (\cref{sec:efficiency-exposure}).
\end{theorybox}

\subsection*{The inverse-beta-decay detection principle}
\label{sec:ibd-principle}

\begin{theorybox}
Every real event this experiment (and this document's inference) is built
from is an inverse-beta-decay (IBD) interaction,
\begin{equation}
  \keyresult{\bar\nu_e + p \rightarrow e^+ + n,}
  \label{eq:intro-ibd-reaction}
\end{equation}
occurring on a free (unbound hydrogen) proton in the GdLS target. IBD is
selected via a real, two-part delayed-coincidence signature, occurring
within the same detector:
\begin{description}
  \item[Prompt signal.] The positron deposits its kinetic energy locally and
    then annihilates with an atomic electron, $e^+ + e^- \to 2\gamma$
    ($2\times 511$~keV), giving a real visible (``prompt'') energy
    \begin{equation}
      E_{\rm prompt} \approx T_{e^+} + 2 m_e = E_\nu - \Delta_{np} + m_e,
      \qquad \Delta_{np} = M_n - M_p,
    \end{equation}
    at leading order (the full order-1/M kinematics are developed in
    \cref{ch:detector-implementation}).
  \item[Delayed signal.] The IBD neutron thermalises (over microseconds)
    and is captured, predominantly on a real gadolinium nucleus
    (chosen for this purpose precisely because of its large capture cross
    section), releasing a real $\sim 8$~MeV gamma-ray cascade $\sim 30\
    \mu\mathrm{s}$ later on average -- well separated in both time and
    energy from the prompt signal and from most backgrounds.
  \item[Selection cuts.] A sequence of real cuts (flasher rejection,
    capture-time window, prompt/delayed energy windows, muon veto,
    multiplicity) isolates the real IBD candidate sample; \cref{ch:results}
    reports how tpeanuts' own real IBD-selection efficiency compares to the
    official one.
\end{description}
This delayed-coincidence tag is what makes IBD, among the handful of
reactions a reactor antineutrino can undergo in a scintillator detector,
overwhelmingly the cleanest and most background-suppressible channel --
exactly why every reactor-antineutrino oscillation experiment to date
(Daya Bay, RENO, Double Chooz, KamLAND) is built around it.
\end{theorybox}

\begin{notebox}
The real coincidence window, multiplicity cuts, and calibration procedures
that turn raw PMT waveforms into a validated IBD candidate list are
themselves a substantial piece of the Collaboration's own analysis
infrastructure; \texttt{tpeanuts} does not reproduce this reconstruction
chain, and instead consumes its real, already-processed output directly
(the official released event-count tables, \cref{ch:data-release}) as its
external, unmodified real-data input.
\end{notebox}

\section{Physics goals and data-taking timeline}

\begin{theorybox}
Real data taking proceeded in three real sub-periods, distinguished by how
many of the 8 ADs were physically installed and operating:
\begin{description}
  \item[6AD period] (December 2011 -- July 2012): 6 ADs online (2 at EH1,
    1 at EH2, 3 at EH3) -- the dataset behind the original 2012 discovery
    \parencite{An2012DayaBayDiscovery}.
  \item[8AD period] (October 2012 onward): all 8 ADs online (2 at each near
    hall, 4 at the far hall) -- the largest and most-used real dataset, and
    the \emph{only} period this document's \texttt{tpeanuts} model fits
    (\cref{ch:detector-implementation}).
  \item[7AD period] (a later real sub-interval): AD11 temporarily removed
    from physical data taking.
\end{description}
The physics goal evolved with the data: the 6AD-period near/far rate
comparison established $\theta_{13}\neq0$
\parencite{An2012DayaBayDiscovery}; with the accumulated 8AD statistics, a
full spectral-shape analysis (comparing the real reconstructed
prompt-energy spectrum bin by bin, not just the total rate) additionally
constrains the mass-squared splitting via the energy-dependent position of
the oscillation minimum, reported as the effective splitting $\Delta
m^2_{ee}$ \parencite{An2023DayaBayPrecision}. \Cref{ch:inference-process}
documents both strategies (rate-only and full-spectral-shape) as
implemented in \texttt{tpeanuts}, plus a real near/far \emph{ratio} fit
that recovers both parameters jointly while remaining robust to the
detector-response simplifications made in this project
(\cref{ch:results}).
\end{theorybox}
